% 2026-11-04. P0 limited source acceptance and exact P1 Coulomb test.
\section{P0 限定源项冻结与 P1 原 Coulomb 磁能的解析关口}
\label{sec:p1-exact-coulomb-gate}

\subsection{P0 的可验证结论及适用域}
本机在 \texttt{21daadf} 上
83 项单元测试通过，P0 两端口各 84 个海水
面/棱/角/远场采样点上，
生产源四阶解析直线与曲角 Gauss 方法
的 $\bm C$、$\bm B$ 相对误差为
$7.28\times10^{-5}$、$4.27\times10^{-4}$，
曲角 24--32 阶高阶独立参考近机器一致。
另一种基于原 BFZI 曲角的物理闭合折线
对 16、32、64 弦段的相对磁场误差
分别约 $3.90\times10^{-4}$、
$9.76\times10^{-5}$、$2.44\times10^{-5}$，
呈近似二阶衰减，且每一折线
$\Div\bm C_{\rm poly}$ 在导线外均为机器零。

因此可以将\emph{给定物理参数和理想闭合单位线源}
的 P0 源积分定义与对照视为限定通过，
不能将之拓展为全参数域、有限线径的自碰撞
或真实铜线近邻电流分布的认证。
固定相量方向、回流路径和材料条件，
并明确人工有限绝缘外球仍非无限导电海水。

\subsection{严格符合电流约束的解析球形原核基准}
在 $\Omega_R=\{\bm x\in\RR^3:|\bm x|<R\}$，
选取三个真实电流向量基：
\[
\bm W_i(\bm x)=\bm e_i\times\bm x,\qquad i=1,2,3.
\]
各模式满足
\[
\Div\bm W_i=0,\qquad
\bm n\cdot\bm W_i|_{\partial\Omega_R}=0,
\qquad
\bm W_i=\Curl\bm F_i,\qquad
\bm F_i=\frac{R^2-|\bm x|^2}{2}\bm e_i,
\quad
\bm F_i|_{\partial\Omega_R}=0 .
\]
这是具有连续物理零法向边界条件和零迹旋度势的
严格电流模式，不是随意选择的常向量试验。

对任一笛卡尔分量为 $\bm W_i$ 的线性函数，
定义原 Coulomb Newton 势
\[
\bm U_i(\bm x)
=\int_{\Omega_R}
\frac{\bm W_i(\bm y)}{|\bm x-\bm y|}\,
\mathrm d\bm y .
\]
利用 $\Delta_{\bm x}\bm U_i=-4\pi\bm W_i$，
球内 $l=1$ 正则径向解
$\bm U_i=c(r)\bm W_i$，
以及球外衰减偶极项的连续匹配，可得
\begin{equation}
\boxed{
\bm U_i(\bm x)
=\left(\frac{2\pi R^2}{3}
-\frac{2\pi|\bm x|^2}{5}\right)
\bm W_i(\bm x),\quad |\bm x|\le R .
}
\label{eq:p1-ball-exact-potential}
\end{equation}
这里不含任何软化、替代球云电流模型或数值
对角自项定义，且边界与标准 $1/r$ 核均为真实积分。

球的各向同性使不同坐标轴模式正交。
由于
\[
\int_{\Omega_R}|\bm W_i|^2\,\mathrm d\bm x
=\frac{8\pi R^5}{15},\qquad
\int_{\Omega_R}|\bm x|^2|\bm W_i|^2
\,\mathrm d\bm x=\frac{8\pi R^7}{21},
\]
代入原核双重积分直接得到
\begin{align}
L_{ij}
&=\frac{\mu_0}{4\pi}
\iint_{\Omega_R\times\Omega_R}
\frac{\bm W_i(\bm x)\cdot\bm W_j(\bm y)}
{|\bm x-\bm y|}
\,\mathrm d\bm y\,\mathrm d\bm x\\
&=\boxed{\frac{16\pi\mu_0 R^7}{315}\,\delta_{ij}}.
\label{eq:p1-ball-analytic-L}
\end{align}
这些模式\emph{没有按安培载流量归一化}，
故其算子双线性型带模式尺寸因子，不得
把该数值误标为真实两端口电感的亨利数值。
本关以完全已知的对角/非对角三模态磁能矩阵
提供\textbf{可证伪的原核数值参考}。

\subsection{相同电流、不同积分轴与局部核分区}
\texttt{examples/bfzi\_p1\_coulomb.py}
使用独立球形 Gauss \texttt{x/y} 积分规则，
直接调用当前未软化原核
\texttt{assemble\_coulomb\_split}：
\begin{enumerate}
\item 固定上述 $\Omega_R,\bm W_i$，
将外积分轴 $x$、内积分轴 $y$ 以
\texttt{base/x\_only/y\_only/both} 四组比较
\emph{同一个解析真值}；
\item 固定节点和电流，改变局部
分区球半径 $a/R=0.25,0.55$，
检查结果对\emph{数学上等价的原核分区}
的敏感度，不改变核 $1/r$；
\item \textbf{在本机首次运行前冻结球基准专用验收阈值：}
用最终 \texttt{both} 规则与式
\eqref{eq:p1-ball-analytic-L} 比较，
矩阵 Frobenius 相对误差、
最大对角元相对误差、非对角相对误差
均须小于 $10^{-3}$；
以解析 $L_{11}$ 归一化的
局部分区 $a/R=0.25,0.55$
双规则差也须小于 $10^{-3}$。
\texttt{base/x\_only/y\_only/both} 的序列仍须人工
分析是否稳定逼近同一真值，
不允许只因最后一个样本偶然落在阈值内就认证。
\item 记录原双重积分三模态矩阵与式
\eqref{eq:p1-ball-analytic-L} 的整体
Frobenius 相对误差、最大对角相对误差、
非对角误差及对称性偏差；
\item 若球基准未达到有证据的稳定性，
先查核奇异区、局部求积与
内外积分轴，而不继续扩大 229 模态字典，
不更换成 Plummer 或实体球云近似。
\end{enumerate}
本机独立入口为
\texttt{python run.py --mode bfzi-p1-coulomb --with-tests}，
独立版本化日志在
\texttt{docs/research/bfzi\_uwpt\_electrothermal/runs/p1/}。
当前\textbf{该 P1 新算例尚未取得本机数值运行证据}，
不能宣布原核求积成功。


\subsection{P1 原核球基准实测失败及严格误差归因}
用户本机 \texttt{ba16d6a} 的第一个 P1 运行
85 项程序单测通过，但\emph{原核解析准确度失败}。
在 $R=1\,{\rm m}$ 的三解析无散度模式上，
当前分区算子 \texttt{assemble\_coulomb\_split}
分别以基础、只加密 $x$、只加密 $y$、
双轴加密的体积规则得到 Frobenius 相对误差
$2.120667\%$、$1.556394\%$、
$1.146122\%$、$0.614486\%$，
最后一项仍高于预定 $0.1\%$ 目标。
固定 $x/y$ 规则，局部分区半径 $a/R$
改为 $0.25,0.55$ 后，
$L_{11}$ 的差值相对解析真值为 $0.4230034\%$。
对称性、非对角元误差均接近机器零，
没有提供数值正确性的充分证据。

目前的 C$^2$ 局部径向窗函数满足
$\chi(s)=1$（$s\le1/2$），
$\chi(s)=0$（$s\ge1$），中间为
\[
\chi(s)=1-10(2s-1)^3+15(2s-1)^4-6(2s-1)^5,
\qquad \frac12<s<1,
\]
故其局部矩可\emph{完全解析}：
\[
\int_0^1 s\chi(s)\,\mathrm ds=\frac27 .
\]
令 $a(\bm x)=
\min(a_{\max},0.6\,\mathrm{dist}(\bm x,\partial\Omega_R))$，
局部球包含于导电球内部，
且 $\bm W_i(\bm y)=
\bm W_i(\bm x)+\bm e_i\times(\bm y-\bm x)$；
后项在对称球上积分为零，于是
\begin{equation}
\boxed{
\int_{|\bm y-\bm x|<a(\bm x)}
\frac{\chi(|\bm y-\bm x|/a)}
{|\bm x-\bm y|}\bm W_i(\bm y)\,\mathrm d\bm y
=\frac{8\pi}{7}\,a(\bm x)^2\,\bm W_i(\bm x).
}
\label{eq:p1-analytic-local-window}
\end{equation}
这一式直接给出\emph{原积分分区}的确切局部值，
无需再加密局部球的径向、极向或方位求积。
已有球内势的解析解
\eqref{eq:p1-ball-exact-potential} 则可直接在
$x$ 规则上计算原算子而不经过任何 $y$ 求积。
两者相减就能将误差定位为：
\begin{enumerate}
\item 解析内势经 $x$ 求积的外积分误差；
\item 局部球采用式
\eqref{eq:p1-analytic-local-window} 精确值后，
$y$ 远区规则相对解析内势的误差；
\item 原程序给出的局部球数值值与上述精确值的差。
\end{enumerate}
其中远区核
$(1-\chi(|\bm x-\bm y|/a(\bm x)))/|\bm x-\bm y|$
虽然在对角线取零，但靠近球壁时
$a(\bm x)\to0$，转折层收缩。
固定全局球 Gauss 求积容易漏掉这类细小空间尺度。
新 \texttt{p1\_decomp} 旨在直接证实或反驳
该误差来源，不能凭直觉认定它已经通过。

\subsection{保留物理原核的全域奇异常量减除}
对任何具有足够正则性的测试场，恒等式
\begin{equation}
\boxed{
\int_\Omega\frac{\bm W(\bm y)}{|\bm x-\bm y|}
\,\mathrm d\bm y
=\bm W(\bm x)\,S_\Omega(\bm x)
+\int_\Omega
\frac{\bm W(\bm y)-\bm W(\bm x)}{|\bm x-\bm y|}
\,\mathrm d\bm y,
\quad
S_\Omega(\bm x)=\int_\Omega
\frac{\mathrm d\bm y}{|\bm x-\bm y|}.
}
\label{eq:p1-global-constant-subtraction}
\end{equation}
当 $\bm W$ 局部 Lipschitz 时，
最后一项在 $y\to x$ 的函数值有界，
但一般仍有依赖方向的极限，不得在
$x/y$ 节点重合处任意赋值。
这是\textbf{未软化的原始 Coulomb 核}，
且并非均匀小球电流云替代模型。

球形几何的标量势已知：
\[
S_{B_R}(\bm x)=2\pi
\left(R^2-\frac{|\bm x|^2}{3}\right).
\]
据此新增
\texttt{green\_coulomb.assemble\_coulomb\_subtracted}，
直接用给定真实几何的独立标量势
和独立 $x/y$ 体积节点，对式
\eqref{eq:p1-global-constant-subtraction}
构成原核磁能矩阵，作为第二种 P1 算法候选。
当前仅能在\emph{球上}把解析
$S_{B_R}$ 当成数值真值。

对双盒排除的海水域 \emph{不可以直接使用}
球形 $S_{B_R}$。
若要使该算法应用于真实
$\Omega=B_R\setminus(\overline B_1\cup\overline B_2)$，
仍需独立构造并验证
$S_\Omega(\bm x)$。
用散度恒等式
$\Div_{\bm y}[(\bm y-\bm x)/|\bm y-\bm x|]
=2/|\bm y-\bm x|$，
可将其形式上写成物理边界的曲面积分：
\[
\boxed{
S_\Omega(\bm x)=\frac12
\int_{\partial\Omega}
\frac{\bm n_{\bm y}\cdot(\bm y-\bm x)}
{|\bm y-\bm x|}\,\mathrm dS_{\bm y}.
}
\]
外球使用外法向，而内盒壁使用\emph{海水域外法向}
（即朝向绝缘盒内部）。
该表面表示本身需要近壁误差控制，
\textbf{在未经证明或核对前不能称双盒算法已通过}。
已新增诊断和原核减除实现，
但这些变更尚未取得用户本机 P1 新实测。


\subsection{原核进一步减除一阶几何矩：非线性解析涡流验算}
仅使用式 \eqref{eq:p1-global-constant-subtraction}
减去常量时，剩余被积核虽然有界，
但对角附近仍具有角向极限。
对有合适 $C^2$ 正则性的电流，
可以进一步将 \emph{完全相同的物理原核}
分解为
\begin{align}
\bm U_W(\bm x)
={}&\bm W(\bm x)\,S_\Omega(\bm x)
+\nabla\bm W(\bm x)\,\bm M_\Omega(\bm x)
\notag\\
&+
\int_\Omega
\frac{
 \bm W(\bm y)-\bm W(\bm x)
 -\nabla\bm W(\bm x)(\bm y-\bm x)}
 {|\bm x-\bm y|}\,\mathrm d\bm y ,
\label{eq:p1-affine-taylor-subtraction}\\
\bm M_\Omega(\bm x)
={}&\int_\Omega\frac{\bm y-\bm x}{|\bm x-\bm y|}
\,\mathrm d\bm y
=\int_{\partial\Omega}
 |\bm x-\bm y|\bm n_{\bm y}\,\mathrm dS_{\bm y}.
\label{eq:p1-geometry-first-moment}
\end{align}
最后一式由
$\nabla_{\bm y}|\bm x-\bm y|=
(\bm y-\bm x)/|\bm y-\bm x|$
及散度定理给出。
余项因一阶 Taylor 减除在对角仅为
$O(|\bm x-\bm y|)$，
从数学上比原始 $1/r$ 核和零阶减除更规则。
\textbf{式 \eqref{eq:p1-affine-taylor-subtraction}
是严格恒等式，不是软化或新的物理模型}。

球内的两个几何矩具有独立的解析值：
\[
S_{B_R}(\bm x)=
2\pi\left(R^2-\frac{|\bm x|^2}{3}\right),\qquad
\boxed{\bm M_{B_R}(\bm x)=
\frac{4\pi}{15}(|\bm x|^2-5R^2)\bm x.}
\]
若仅用一阶电流 $\bm W_i=\bm e_i\times\bm x$，
Taylor 余项\emph{严格为零}，因而该例只能
验证 $\bm S,\bm M$ 和球外积分，
\textbf{不能}据此认证余项的实际求积。

为防止上述制造解过于平凡，
新增一个\emph{非线性、严格无散度、
球壁零法向} 的第二级真值模式：
\[
\bm W_i^{(3)}(\bm x)
=|\bm x|^2(\bm e_i\times\bm x),\qquad
\bm F_i^{(3)}(\bm x)
=\frac{R^4-|\bm x|^4}{4}\bm e_i ,
\quad \Curl\bm F_i^{(3)}=\bm W_i^{(3)} .
\]
从 $\Delta\bm U_i^{(3)}=-4\pi\bm W_i^{(3)}$
及球外 $l=1$ 匹配得到
\begin{equation}
\boxed{\bm U_i^{(3)}(\bm x)
=\left(\frac{\pi R^4}{3}
-\frac{\pi|\bm x|^4}{7}\right)
(\bm e_i\times\bm x),\quad
L_{ij}^{(3)}=\frac{16\pi\mu_0 R^{11}}{693}\delta_{ij}.}
\label{eq:p1-cubic-vortex-exact}
\end{equation}
这个模式的 Taylor 余项一般\emph{不为零}，
可以真正检验式
\eqref{eq:p1-affine-taylor-subtraction}
的残余 $y$ 积分。
独立短程球体积数值计算在固定 480 个 $x$、
756 个 $y$ 节点下发现：
仅减常量的非线性电流磁矩阵相对误差约
$2.453\times10^{-3}$，
进一步减除一阶项后约
$1.596\times10^{-4}$，
而解析内势仅在 $x$ 节点积分类似误差约
$5.91\times10^{-6}$。
这些是独立短程校验数值，
\textbf{新增 GitHub P1 自动日志与单测尚待本机复跑}，
不能将其等同于完整 BFZI 教师验收。

生产代码的\texttt{assemble\_coulomb\_subtracted}
现在支持配对参数
\texttt{gradient}、\texttt{first\_moment}，
并坚持输入 \emph{已独立可信的}
实际几何标量势和一阶矩。
对有两个旋转绝缘盒孔洞的海水域，
不能直接代入球公式：
\textbf{真实双盒 $S_\Omega,\bm M_\Omega$
和临近盒壁的相应积分必须分别认证}。
这一步没有改动原 BFZI 端口源、
$R/K$ 或神经网络教师训练。

\subsection{与真实 BFZI 双盒积分之间的限制}
即使式 \eqref{eq:p1-ball-analytic-L} 在球上成功，
它只证明在此解析无散度电流空间和球形边界
下的原 Coulomb 核积分经过了一个独立基准。
BFZI 的真实海水域是
两只独立 SE(3) 旋转绝缘盒的外部，
存在非光滑棱角和多体近邻相互作用，
另需对\emph{同一固定物理电流场} 的对角块、
靠近盒面的局部积分与远区子块单独认证。
涡流体积积分中的奇异电感自项与
分离内外积分轴的困难见
Passarotto--Pitassi--Specogna~\cite{passarotto2022}；
式 \eqref{eq:g1-review-subtraction} 的解析
奇异减除仍是后续优先途径，
不可以用球形算例代替所有边界角度验证。
